% skeleton-paper-talk.tex -- sv-slides paper-talk on Arkhangelsky et al. SDID
% Numbers copied from skills/study/evals/.../quickread.md (Table 1, arXiv:1812.09970).
% Theme: simpleplus-cjk. Compile:
%   skills/sv-slides/scripts/build.sh --theme simpleplus-cjk templates/skeleton-paper-talk.tex
\documentclass[aspectratio=169,10pt]{beamer}
\usepackage{groupbeamer-simpleplus}

\title{Synthetic Difference in Differences}
\subtitle{Critical paper talk (group meeting)}
\author{Arkhangelsky, Athey, Hirshberg, Imbens, Wager}
\institute{arXiv:1812.09970 (July 2021 version)}
\date{Presenter / 2026-10-17}

\begin{document}

\begin{frame}
  \titlepage
  \note{State the one-line claim: weighted TWFE with unit and time weights.}
\end{frame}

\begin{frame}{The aim is causal: ATT for a block-treated panel}
  \begin{equation*}
    \tau = \frac{1}{N_{\mathrm{tr}}T_{\mathrm{post}}}
    \sum_{i \in \mathrm{tr}}\sum_{t \in \mathrm{post}}
    \bigl(Y_{it}(1)-Y_{it}(0)\bigr)
  \end{equation*}
  \begin{itemize}
    \item Unit, time, intervention: states, years, a single adoption date.
    \item Presenter's reading of Eq.~4.2 (block assignment). Staggered
      adoption is backup, not this talk.
  \end{itemize}
  \note{If someone asks about staggered adoption, defer to the appendix.}
\end{frame}

\begin{frame}{Paper card (presenter's filling)}
  \vspace{4pt}
  \begin{tabular}{ll}
    \toprule
    Field & This paper \\
    \midrule
    Estimand & Block ATT (Eq.~4.2); labelled if vague \\
    Identification & Low-rank \(L\) plus \(E[E\mid W,L]=0\) \\
    Key assumption & Assignment may depend on \(L\), not on \(E\) \\
    Main result & CA: SDID \(-15.6\) (se \(8.4\)), Table 1 \\
    Biggest threat & Placebo designs are zero-effect \\
    \bottomrule
  \end{tabular}
  \note{Unstated bits in the card are flagged in the quickread, not invented.}
\end{frame}

\begin{frame}{Identification: knowledge vs convenience}
  \begin{itemize}
    \item A1 outcome \(Y=L+W\circ\tau+E\), \(L\) low rank (knowledge).
    \item A2 \(E[E\mid W,L]=0\) (knowledge; untestable).
    \item A3 unit-independence of \(E\), homoskedasticity (convenience
      for placebo se).
  \end{itemize}
  \begin{equation*}
    E[E \mid W, L] = 0
  \end{equation*}
  \begin{block}{What breaks}
    If assignment depends on \(E\), \(\tau\) is not identified; a better
    estimator does not close that gap.
  \end{block}
  \note{Rank A2 as the strong one. Do not upgrade A3 into identification.}
\end{frame}

\begin{frame}{Method: unit weights, time weights, then TWFE}
  \begin{equation*}
    (\hat\tau^{\mathrm{sdid}},\hat\mu,\hat\alpha,\hat\beta)
    = \arg\min \sum_{i,t}
    \omega_i\lambda_t \bigl(Y_{it}-\mu-\alpha_i-\beta_t-W_{it}\tau\bigr)^{2}
  \end{equation*}
  \begin{itemize}
    \item \(\omega\): pre-treatment trends parallel, intercept allowed.
    \item \(\lambda\): pre periods that predict the post mean.
  \end{itemize}
  \note{Eq.~1.1 in the paper. DID is the same regression with flat weights.}
\end{frame}

\begin{frame}{DAG sketch: assignment may use \(L\), not \(E\)}
  \begin{tikzpicture}[
      node distance=1.4cm, auto,
      obs/.style={circle, draw, minimum size=7mm},
      cf/.style={circle, draw, dashed, minimum size=7mm}
    ]
    \node[obs] (L) {\(L\)};
    \node[obs, right=of L] (W) {\(W\)};
    \node[obs, below=of W] (Y) {\(Y\)};
    \node[cf, left=of Y] (E) {\(E\)};
    \draw[->] (L) -- (W);
    \draw[->] (L) -- (Y);
    \draw[->] (W) -- (Y);
    \draw[->, dashed] (E) -- (Y);
  \end{tikzpicture}
  \note{Solid observed factors. Dashed error. No arrow from E to W under A2.}
\end{frame}

\begin{frame}{California pack-sales estimate (Table 1)}
  \vspace{4pt}
  \begin{tabular}{lrr}
    \toprule
    Estimator & \(\hat\tau\) & se (placebo) \\
    \midrule
    SDID & \(-15.6\) & \(8.4\) \\
    SC   & \(-19.6\) & \(9.9\) \\
    DID  & \(-27.3\) & \(17.7\) \\
    \bottomrule
  \end{tabular}
  \begin{itemize}
    \item Units: packs per person per year, 1970--2000, 38 controls.
  \end{itemize}
  \note{Read the numbers. Do not say SDID caused a smaller effect.}
\end{frame}

\begin{frame}{Causal reading holds only under A1--A2}
  \begin{itemize}
    \item Under A1--A2, SDID estimates the block ATT, not DID's ATT
      under parallel trends.
    \item Sensitivity: none reported for assignment-on-\(E\).
    \item Placebo RMSE (CPS baseline): SDID \(0.028\), DID \(0.049\)
      (Table 2, 1000 draws). Zero-effect design.
  \end{itemize}
  \begin{equation*}
    \widehat{\mathrm{RMSE}}_{\mathrm{CPS}}(\mathrm{SDID}) = 0.028
  \end{equation*}
  \note{Hernan split: numbers were the previous slide; this is the reading.}
\end{frame}

\begin{frame}{Link to our project (fill or TODO)}
  % gb:text-ok project link is prose by nature
  \begin{itemize}
    \item If we have few treated units and doubtful parallel trends,
      SDID is a robustness check beside DID and SC.
    \item Confirm a common adoption date before using Eq.~4.2.
    \item Project RQ: TODO.
  \end{itemize}
  \note{Do not claim our data fit SDID until STATUS says so.}
\end{frame}

\begin{frame}{Reproducibility: violate A2 in a placebo}
  \begin{equation*}
    W_i = \mathbf{1}\{E_{i1}>\mathrm{median}(E_{\cdot 1})\}
  \end{equation*}
  \begin{itemize}
    \item Simulate assignment that depends on \(E\); report bias of SDID
      vs DID vs SC (TODO: not run in this skeleton).
    \item Optional 60-min segment: replicate Table 1 side by side.
  \end{itemize}
  \note{This is a plan, not a result. Keep TODO visible.}
\end{frame}

\begin{frame}{What would change my mind}
  \begin{itemize}
    \item A non-zero-effect simulation where SDID bias exceeds DID.
    \item Formal conditions for staggered adoption (appendix §8 is a
      recipe, not a theorem in the quickread).
  \end{itemize}
  \begin{equation*}
    \hat\tau^{\mathrm{sdid}} - \hat\tau^{\mathrm{did}}
    \quad\text{is not a test of } E[E\mid W,L]=0.
  \end{equation*}
  \begin{block}{Limit}
    Theory in the talk is block assignment only.
  \end{block}
  \note{Invite the group to name another killer check.}
\end{frame}

\begin{frame}{Discussion: should we treat SC/SDID mismatch as a warning?}
  % gb:text-ok closing discussion question
  \begin{itemize}
    \item Paper: when SC and SDID disagree, SDID is usually more
      credible (§2.2). Quickread tag: conjectured.
    \item California has no ground truth, so a mismatch alone does not
      tell us which estimator is wrong.
  \end{itemize}
  \note{End on the question, not a slogan.}
\end{frame}

\begin{frame}{References}
  \vspace{4pt}
  \begin{tabular}{ll}
    \toprule
    Source & Use \\
    \midrule
    arXiv:1812.09970 & Tables 1--2 (quickread card) \\
    arXiv:2305.18793 & Ding notation defaults \\
    \bottomrule
  \end{tabular}
\end{frame}

\begin{frame}{Thank you}
  \begin{equation*}
    \tau \;=\; E[Y(1)-Y(0)\mid \text{block assignment}]
  \end{equation*}
\end{frame}

\appendix
\begin{frame}{Backup: staggered adoption is not this estimand}
  \begin{equation*}
    \hat\tau = \sum_{g} \pi_g \hat\tau_g^{\mathrm{sdid}}
  \end{equation*}
  \note{Appendix §8 recipe. No theorem in the card. Hyperlink from slide 2.}
\end{frame}

\end{document}
